
\bta{2006甲B}

\begin{enumerate}
	%\renewcommand{\labelenumi}{\arabic{enumi}.}
	% A(\Alph) a(\alph) I(\Roman) i(\roman) 1(\arabic)
	%设定全局标号series=example	%引用全局变量resume=example
	%[topsep=-0.3em,parsep=-0.3em,itemsep=-0.3em,partopsep=-0.3em]
	%可使用leftmargin调整列表环境左边的空白长度 [leftmargin=0em]
	\item
（$30'$）已知谐振子处于第$n$个定态中，试导出算符$\hat{x},\hat{p},(\hat{x})^2,(\hat{p})^2$的平均值及不确定度$\Delta x,\Delta p$，并求出$\Delta x\cdot\Delta p$值。


\banswer{
	令
	\[
		x=\sqrt{\frac{\hbar}{2m\omega}}(a+a^\dagger),
		\qquad
		p=i\sqrt{\frac{m\hbar\omega}{2}}(a^\dagger-a).
	\]
	在第 $n$ 个定态中，$\langle a\rangle=\langle a^\dagger\rangle=0$，并利用
	$\langle a^\dagger a\rangle=n$、$\langle aa^\dagger\rangle=n+1$，得
	\[
		\langle x\rangle=\langle p\rangle=0,
	\]
	\[
		\langle x^2\rangle=\frac{\hbar}{m\omega}\left(n+\frac12\right),
		\qquad
		\langle p^2\rangle=m\hbar\omega\left(n+\frac12\right).
	\]
	因此
	\[
		\boxed{\Delta x=\sqrt{\frac{\hbar}{m\omega}\left(n+\frac12\right)}},
		\qquad
		\boxed{\Delta p=\sqrt{m\hbar\omega\left(n+\frac12\right)}},
	\]
	以及
	\[
		\boxed{\Delta x\,\Delta p=\hbar\left(n+\frac12\right)}.
	\]
}


\item 
（$30'$）设$\hat{\mu}$为幺正算符，若存在两个厄米算符$\hat{A}$和$\hat{B}$使
$\hat{\mu}=\hat{A}+i\hat{B}$。试证：
\begin{enumerate}
	%\renewcommand{\labelenumi}{\arabic{enumi}.}
	% A(\Alph) a(\alph) I(\Roman) i(\roman) 1(\arabic)
	%设定全局标号series=example	%引用全局变量resume=example
	%[topsep=-0.3em,parsep=-0.3em,itemsep=-0.3em,partopsep=-0.3em]
	%可使用leftmargin调整列表环境左边的空白长度 [leftmargin=0em]
	\item
	$\hat{A}^2+\hat{B}^2=1$，且$[\hat{A},\hat{B}]=0$；

\item 
进一步再证明$\hat{\mu}$可表示成$\hat{\mu}=e^{i\hat{H}}$，$\hat{H}$为厄米算符。
\end{enumerate}

\banswer{
	\begin{enumerate}
		\item 因 $\mu$ 幺正，$\mu^\dagger\mu=\mu\mu^\dagger=1$。又 $\mu^\dagger=A-iB$，所以
		\[
			A^2+B^2+i[A,B]=1,\qquad
			A^2+B^2-i[A,B]=1.
		\]
		两式相加、相减即得
		\[
			\boxed{A^2+B^2=1,\qquad[A,B]=0}.
		\]

		\item 对易的厄米算符 $A,B$ 可同时对角化。在任一共同本征态上，其实本征值 $a,b$ 满足 $a^2+b^2=1$，故可选实数 $h$ 使 $a=\cos h$、$b=\sin h$。用谱分解把这些 $h$ 作为本征值定义厄米算符 $H$，便有
		\[
			A=\cos H,\qquad B=\sin H.
		\]
		于是
		\[
			\boxed{\mu=A+iB=\cos H+i\sin H=e^{iH}},
		\]
		其中 $H$ 的本征值可按模 $2\pi$ 选取，故 $H$ 一般不唯一。
	\end{enumerate}
}


\item 
（$30'$）一个质量为$m$的粒子被限制在$0\le x\le a$的一维无穷深势阱中。初始时刻其归一化波函数为$\psi(x,0)=\sqrt{\frac{8}{5a}}(1+\cos\frac{\pi x}{a})\sin\frac{\pi x}{a}$，求：
\begin{enumerate}
	%\renewcommand{\labelenumi}{\arabic{enumi}.}
	% A(\Alph) a(\alph) I(\Roman) i(\roman) 1(\arabic)
	%设定全局标号series=example	%引用全局变量resume=example
	%[topsep=-0.3em,parsep=-0.3em,itemsep=-0.3em,partopsep=-0.3em]
	%可使用leftmargin调整列表环境左边的空白长度 [leftmargin=0em]
	\item
$t>0$时粒子的状态波函数；

\item 
在$t=0$与$t>0$时在势阱的左半部发现粒子的概率是多少？
\end{enumerate}


\banswer{
	令 $\phi_n(x)=\sqrt{2/a}\sin(n\pi x/a)$，$E_n=n^2\pi^2\hbar^2/(2ma^2)$。由
	\[
		(1+\cos\frac{\pi x}{a})\sin\frac{\pi x}{a}
		=\sin\frac{\pi x}{a}+\frac12\sin\frac{2\pi x}{a},
	\]
	题给初态为
	\[
		\psi(x,0)=\frac2{\sqrt5}\phi_1(x)+\frac1{\sqrt5}\phi_2(x).
	\]
	\begin{enumerate}
		\item
		\[
			\boxed{\psi(x,t)=\frac2{\sqrt5}e^{-iE_1t/\hbar}\phi_1(x)
			+\frac1{\sqrt5}e^{-iE_2t/\hbar}\phi_2(x)}.
		\]

		\item 展开模平方，
		\[
		 |\psi|^2=\frac45\phi_1^2+\frac15\phi_2^2
		 +\frac45\phi_1\phi_2\cos\frac{(E_2-E_1)t}{\hbar}.
		\]
		每个定态在左半阱的概率为 $1/2$。干涉项的积分为
		\[
		 \int_0^{a/2}\phi_1\phi_2dx
		 =\frac1a\int_0^{a/2}\left(\cos\frac{\pi x}{a}
		                         -\cos\frac{3\pi x}{a}\right)dx
		 =\frac4{3\pi}.
		\]
		代入得
		\[
			\boxed{P_L(t)=\frac12+\frac{16}{15\pi}
			\cos\frac{(E_2-E_1)t}{\hbar}}.
		\]
		特别地，$P_L(0)=1/2+16/(15\pi)$。
	\end{enumerate}
}


\newpage
\item 
（$30'$）粒子在一维无限深方势阱中运动，受到微扰$H'=\frac{V_0}{a}(a-|2x-a|)$的作用。求第$n$个能级的一级近似，并分析所的结果的适用条件。

\banswer{
	取无扰本征函数 $\phi_n=\sqrt{2/a}\sin(n\pi x/a)$。微扰关于 $x=a/2$ 对称，故
	\begin{align*}
		E_n^{(1)}
		&=\frac{8V_0}{a^2}\int_0^{a/2}x\sin^2\frac{n\pi x}{a}\,dx\\
		&=\frac{V_0}{2}+\frac{V_0[1-(-1)^n]}{n^2\pi^2}.
	\end{align*}
	所以第 $n$ 个能级的一级近似为
	\[
		\boxed{E_n\simeq\frac{n^2\pi^2\hbar^2}{2ma^2}
		+\frac{V_0}{2}+\frac{V_0[1-(-1)^n]}{n^2\pi^2}}.
	\]
	即偶数 $n$ 的修正为 $V_0/2$，奇数 $n$ 的修正为 $V_0/2+2V_0/(n^2\pi^2)$。

	微扰展开要求不同无扰能级间的混合很小：对所有重要的 $m\ne n$，
	\[
		|\langle m|H'|n\rangle|\ll|E_m^{(0)}-E_n^{(0)}|.
	\]
	一个充分的低能级尺度条件是 $|V_0|\ll\pi^2\hbar^2/(2ma^2)$；对高 $n$，可与相邻能级间隔比较而适当放宽。
}


\item 
一个质量为$m$ 的粒子被限制在 $r = a $和 $r = b$ 的两个不可穿透的同心球面之间运动,不存在其它势。求粒子的基态能量和归一化的波函数。

\banswer{
假定 $0<a<b$。令 $\psi=R_l(r)Y_{lm}(\theta,\varphi)$、$u_l=rR_l$，则
\[
 \left[-\frac{\hbar^2}{2m}\frac{d^2}{dr^2}
       +\frac{\hbar^2l(l+1)}{2mr^2}\right]u_l=Eu_l,
 \qquad u_l(a)=u_l(b)=0.
\]
所有 $l$ 的径向定义区间及端点条件相同，而离心势非负，因此最低能量在 $l=0$ 通道取得。
此时 $u''+k^2u=0$。由左端条件写成 $u=C\sin[k(r-a)]$，右端条件给出
$k(b-a)=n\pi$（$n=1,2,\ldots$）。基态取 $n=1$，故
\[
 \boxed{E_0=\frac{\pi^2\hbar^2}{2m(b-a)^2}}.
\]
球谐函数取归一化的 $Y_{00}=1/\sqrt{4\pi}$，三维归一化就等价于
\[
 1=\int_a^b|R_0|^2r^2dr=|C|^2\int_a^b
       \sin^2\frac{\pi(r-a)}{b-a}dr=|C|^2\frac{b-a}{2}.
\]
所以全波函数为
\[
 \boxed{\psi_0(r,\theta,\varphi)=\begin{cases}
 \dfrac{1}{\sqrt{2\pi(b-a)}}\dfrac{\sin[\pi(r-a)/(b-a)]}{r},&a<r<b,\\[6pt]
 0,&\text{其他区域}.
 \end{cases}}
\]
其中 $1/r$ 来自 $R=u/r$，不能把径向函数直接当作一维正弦函数。
}

	
\end{enumerate}
